\section{Conclusiones}

\subsection{Resultados relevantes}
El desarrollo de SAGAI demostró que es completamente viable integrar Inteligencia Artificial en procesos administrativos y legales. Se logró construir un sistema web funcional capaz de extraer texto de documentos digitalizados y conectarlo con la API de Claude para generar borradores de acuerdos judiciales. El mayor logro de este proyecto fue agilizar la redacción técnica sin quitarle el control al usuario, garantizando que el personal del juzgado siempre tenga la revisión y validación final antes de aprobar cualquier documento.

\subsection{Experiencia personal y profesional}
Hacer este proyecto fue un gran paso para mí. La verdad, al principio yo quería enfocarme en el backend, pero se me hizo muy difícil por el momento, así que decidí darle una oportunidad al frontend. ¡Y fue una gran decisión! Armar SAGAI me dejó meterme de lleno a crear las pantallas con HTML, CSS y JavaScript, y descubrí que me gusta muchísimo diseñar la experiencia del usuario y hacer que todo funcione fluido. Además, lograr conectar mi interfaz con una inteligencia artificial me quitó el miedo a trabajar con APIs y me dio muchísima confianza en lo que puedo llegar a programar.

\subsection{Sugerencias de mejora}
Aunque el sistema cumple con sus objetivos principales, hay bastante margen para seguir escalándolo. Como sugerencias de mejora a futuro, sería ideal ampliar la variedad de promociones y tipos de juicios que la IA está entrenada para redactar. Además, se podría implementar un sistema de roles y permisos más robusto en la base de datos, diferenciando los accesos entre Jueces, Secretarios y Actuarios. Por último, integrar un motor de OCR más avanzado que sea capaz de leer documentos antiguos o con sellos borrosos aumentaría enormemente la utilidad de la plataforma.